\documentclass[11pt]{article}
\usepackage[a4paper,margin=1in]{geometry}
\usepackage{fontspec}
\usepackage[boldfont=false]{xeCJK}
\setCJKmainfont{FandolSong-Regular.otf}
\usepackage{amsmath}
\usepackage{amssymb}
\usepackage{array}
\usepackage{booktabs}
\usepackage{graphicx}
\usepackage{adjustbox}
\usepackage{enumitem}
\setlist[itemize]{topsep=2pt,partopsep=0pt,parsep=0pt,itemsep=2pt,leftmargin=1.4em}
\setlist[enumerate]{topsep=2pt,partopsep=0pt,parsep=0pt,itemsep=2pt,leftmargin=1.6em}
\usepackage{parskip}
\usepackage{fvextra}
\fvset{breaklines=true,breakanywhere=true,fontsize=\small}
\setlength{\parindent}{0pt}

\begin{document}

\begin{center}
  {\LARGE\bfseries Classification, biodiversity and conservation}\\[0.35em]
  {\large A Level Biology}
\end{center}
\vspace{0.6em}

\section*{Classifying living things}

A \textbf{species} 物种 can be defined in more than one way:

\begin{itemize}
\item the \textbf{biological} species concept — a group whose members can breed together to produce fertile offspring.

\item the \textbf{morphological} species concept — a group whose members look alike.

\item the \textbf{ecological} species concept — a group that fills the same role in its surroundings.

\end{itemize}

\subsection*{The three domains}

The largest groups in \textbf{classification} 分类 are three \textbf{domains} 域:

\begin{itemize}
\item \textbf{Archaea} 古菌 and \textbf{Bacteria} 细菌 — both are \textbf{prokaryotes} 原核生物 (no nucleus). They look similar but differ in their membrane lipids, their ribosomal RNA, and the make-up of their \textbf{cell walls} 细胞壁.

\item \textbf{Eukarya} 真核生物 — all organisms whose cells have a nucleus.

\end{itemize}

\subsection*{The taxonomic hierarchy}

Inside the Eukarya domain, organisms are sorted into a \textbf{taxonomic hierarchy} 分类层级 — a set of smaller and smaller groups: \textbf{kingdom} 界, \textbf{phylum} 门, \textbf{class} 纲, \textbf{order} 目, \textbf{family} 科, \textbf{genus} 属 and species.

The Eukarya are split into four kingdoms:

\begin{itemize}
\item \textbf{Protoctista} 原生生物界 — mostly single-celled (such as \emph{Amoeba}).

\item \textbf{Fungi} 真菌界 — feed by absorbing food; have cell walls of chitin.

\item \textbf{Plantae} 植物界 — make their own food by photosynthesis.

\item \textbf{Animalia} 动物界 — feed on other organisms.

\end{itemize}

Viruses are not placed in these groups. They are classified by the type of \textbf{nucleic acid} 核酸 they contain (DNA or RNA) and whether it is \textbf{single-stranded} 单链 or \textbf{double-stranded} 双链.

\section*{Biodiversity}

An \textbf{ecosystem} 生态系统 is all the living things in an area together with their non-living surroundings. A \textbf{niche} 生态位 is the particular role and place of a species within it.

\textbf{Biodiversity} 生物多样性 can be measured at three levels:

\begin{itemize}
\item the number and range of different ecosystems and \textbf{habitats} 栖息地.

\item the number of species and their relative \textbf{abundance} 丰度 (how common each one is).

\item the genetic variation within each species.

\end{itemize}

\subsection*{Sampling an area}

You cannot count every organism, so you take samples. \textbf{Random sampling} 随机取样 (choosing positions by chance) avoids bias and gives a fair picture. Useful methods are:

\begin{itemize}
\item \textbf{quadrats} 样方 — square frames placed to count or estimate the species inside them.

\item \textbf{transects} 样带 — counting along a line across the area (a line transect records what touches the line; a belt transect counts within a strip).

\item \textbf{mark-release-recapture} 标志重捕法 — for moving animals: catch, mark and release some, then later see what fraction of a new catch is marked (the Lincoln index).

\end{itemize}

To link the spread of a species to \textbf{biotic factors} 生物因素 (living) or \textbf{abiotic factors} 非生物因素 (non-living), you can use Spearman's rank or Pearson's \textbf{correlation} 相关性. To measure the diversity of an area as a single number, you use \textbf{Simpson's index of diversity} 辛普森多样性指数: a higher value means more diverse and usually more stable.

\section*{Conservation}

A species can become \textbf{extinct} 灭绝 (die out completely) because of \textbf{climate change} 气候变化, competition (often from new species), hunting by humans, or the damage and loss of its habitats.

We try to protect biodiversity because other species give us food, medicines and materials, help keep ecosystems stable, and have value in themselves.

\subsection*{Ways to conserve species}

\begin{itemize}
\item \textbf{zoos} 动物园 and \textbf{botanic gardens} 植物园 keep and breed \textbf{endangered species} 濒危物种.

\item protected \textbf{conservation} 保护 areas, such as national parks and marine parks, keep habitats safe.

\item 'frozen zoos' store frozen cells, eggs and sperm, and \textbf{seed banks} 种子库 store seeds for the future.

\end{itemize}

For rare mammals, \textbf{assisted reproduction} can boost numbers: \textbf{in vitro fertilisation} 体外受精 (IVF), \textbf{embryo transfer} 胚胎移植, and \textbf{surrogacy} 代孕 (another female carries the young).

Conservationists also control \textbf{invasive species} 入侵物种, which are brought in from elsewhere and out-compete native species.

Two organisations help worldwide: the IUCN, which lists how threatened each species is, and CITES, which controls the international trade in endangered animals and plants.


\end{document}
